\scanpage{35}

\[
 N^-=\operatorname{Transp}(U_1,U_0)\cap\operatorname{Transp}(U_0,U_1)\tag*{(30)}
\]

Soient maintenant $\beta\in L_0$, $\gamma\in L_1$ avec

\[
\left\{\begin{aligned}
 &\beta\gamma=-1\quad\text{i.e. }\gamma=-\beta^{-1},\ \text{i.e. }\beta=-\gamma^{-1}\\
 &\text{i.e. }\ e_1(\gamma)=\underbrace{\sigma_0(\beta)(e_0(\beta))}_{e_1(-\beta^{-1})},
\end{aligned}\right.\tag*{(31)}
\]

je dis qu'on a alors

\[
 e_1(\gamma)e_0(\beta)e_1(\gamma)=\sigma_0(\beta)\qquad(\text{pour }\underline{\beta\gamma=-1})\tag*{(32)}
\]

formule qui s'écrit aussi, puisque $e_1(\gamma)=\sigma_0(\beta)e_0(\beta)\sigma_0(\beta)^{-1}$

\[
 \underbrace{\sigma_0(\beta)e_0(\beta)}\,\underbrace{\sigma_0(\beta)^{-1}e_0(\beta)}\,\underbrace{\sigma_0(\beta)e_0(\beta)}\,\sigma_0(\beta^{\deleted{\unsure{-1}}})^{-1}=\sigma_0(\beta)
 \tag*{\deleted{\unsure{(33)}}}
\]

soit encore, en écrivant $\sigma_0(\beta)^{-1}=\unsure{\omega|w}\,\sigma_0(\beta)$ ($\omega$ central, $\omega^2=1$)

\[
 (\sigma_0(\beta)e_0(\beta))^3\deleted{\unsure{\omega}}=\deleted{\unsure{\omega}}1
\]

soit encore

\[
 \rho(\beta)^3=1,\quad\text{où }\rho(\beta)=\sigma_0(\beta)e_0(\beta),\tag*{(33)}
\]

soit

\[
 \rho(\beta)=\bigl(\begin{smallmatrix}0&\beta\\-\beta^{-1}&-1\end{smallmatrix}\bigr)\deleted{\unsure{\text{est}}}
 \overset{\deleted{\text{soit }\rho(\beta}}{\underset{\substack{\text{si }L_0=L_1=\Lambda\\\text{donc }\beta\in\Lambda^*}}{=}}
 u(\beta)(\rho')\deleted{\unsure{u}}\quad\rho'=\rho(1)=\bigl(\begin{smallmatrix}0&1\\-1&-1\end{smallmatrix}\bigr)\tag*{(34)}
\]

la vérification est immédiate.

Supposons maintenant $\beta\in L_0$, $\gamma\in L_1$ quelconques, et calculons $e_1(\gamma)e_0(\beta)e_1(\gamma)$, on trouve

\[
 e_1(\gamma)e_0(\beta)(e_1(\gamma))=\bigl(\begin{smallmatrix}1+\beta\gamma&\beta\\\gamma(2+\beta\gamma)&1+\beta\gamma\end{smallmatrix}\bigr)\tag*{\unsure{(36)|(35)}}
\]

% slanted
\margin{faire le lien avec la formule pour $e_1(\deleted{\unsure{\gamma}}\beta)e_0(\deleted{\unsure{\beta}}\alpha')$ comme $e_0(x)u(t)e_1(y)$ pour $1+\deleted{\unsure{\alpha}}\beta\alpha'$ \unsure{invt.}}

qui prouve

\[
\begin{aligned}
 &\text{Si }\beta\in L_0,\gamma\in L_1,\text{ les conditions suivantes sont équivalentes}\\[-1pt]
 &\left\{\begin{aligned}
 \text{a)}\quad&\beta\gamma=-1\\
 \text{b)}\quad&e_1(\gamma)e_0(\beta)e_1(\gamma)\in\operatorname{Transp.}(U_0,U_1)\\
 \text{c)}\quad&e_1(\gamma)e_0(\beta)e_1(\gamma)\in\operatorname{Transp.}(U_1,U_0)\\
 \text{d)}\quad&e_1(\gamma)e_0(\beta)e_1(\gamma)\in N^-\qquad
 \text{e)}\quad e_1(\gamma)e_0(\beta)e_1(\gamma)=\sigma_0(\beta)
 \end{aligned}\right.
\end{aligned}\tag*{(37)}
\]

% archive sheet number: 34
